\documentclass[11pt,a4paper]{article}

\usepackage{amsmath,amssymb,amsthm}
\usepackage[margin=3cm]{geometry}

\newtheorem{theorem}{Theorem}[section]
\newtheorem{lemma}[theorem]{Lemma}
\newtheorem{proposition}[theorem]{Proposition}
\newtheorem{corollary}[theorem]{Corollary}
\theoremstyle{definition}
\newtheorem{definition}[theorem]{Definition}
\newtheorem{remark}[theorem]{Remark}

\newcommand{\Next}{\mathop{\bigcirc}}
\newcommand{\Always}{\mathop{\square}}
\newcommand{\Event}{\mathop{\lozenge}}
\newcommand{\Sub}{\mathrm{Sub}}
\newcommand{\St}{\mathrm{St}}
\newcommand{\Ust}{\mathrm{U}}
\newcommand{\proves}{\vdash_{\mathcal{L}}}
\newcommand{\Reach}{\mathrm{Reach}}
\newcommand{\outs}{\mathrm{out}}

\title{A finite fair-graph certificate\\
{\large and the SCC decision procedure for LTL, made proof-producing}}
\author{}
\date{}

\begin{document}
\maketitle

\begin{abstract}
The standard automata-theoretic decision procedure for LTL satisfiability
decomposes the state graph into strongly connected components: a non-trivial
component is accepting when it meets every fairness set, and a fair path exists
iff an accepting component is reachable.  Its correctness rests on the
assertion that an infinite path in a finite graph eventually remains within one
component, which is infinitary and, in the negative direction, classical.  We
show that the assertion is not needed if the procedure is asked to return a
certificate rather than a bit.  The negative answer is taken to be a
\emph{well-founded record of a removal order} --- at each step either a vertex
with no out-edge, or a successor-closed block missing a fairness set --- and
this record both refutes the existence of a fair cycle, by two inductions on a
finite closed trail, and serves as the induction principle from which Hilbert
proofs are read off.  The resulting statement is a theorem of graph theory with
no formula in it; the temporal content enters only through the choice of
fairness predicates.  Everything has been checked in Agda, and the resulting
procedure is compared, by checked equation, with the one of the companion
note~\cite{m4}.
\end{abstract}

\section*{Introduction}

The companion note~\cite{m4} decides $\proves\varphi$ by an elimination
procedure and reads a Hilbert proof off its failure.  The present note does the
same for the other standard presentation, the one that decomposes the graph
into strongly connected components; the two are then in the same development
and can be compared by computation rather than by assertion.

Three things came out of doing this that are not visible in either informal
account.

First, the graph argument \emph{factors}.  Neither the elimination of~\cite{m4}
nor the usual SCC presentation separates the graph theory from the logic: sets
of states, reachability, components and lassos are all described in terms of
subformulas.  They need not be.  Section~\ref{sec:graph} states and proves a
theorem about a finite graph, finitely many decidable vertex predicates and a
vertex, in which no formula occurs; Sections~\ref{sec:lasso}
and~\ref{sec:extract} are the translation layer.  The choice of predicate is
where the logic goes, and it is one line (Definition~\ref{def:evpred}).

Second, the ``key insight'' is a red herring in the direction it is usually
stated.  What is needed for correctness is not that an infinite path stays in a
component, but the contrapositive statement about a \emph{finite} cycle, and
that is two inductions on a closed trail of the length it happens to have
(Theorem~\ref{thm:nocirc}).  No infinite path is constructed and nothing is
coinductive.

Third, the SCC form gives a genuine improvement on one point and a genuine loss
on another, and both can be measured.  The improvement: in a component that is
bad for an eventuality, \emph{every} state has that eventuality pending, so a
single use of (Ind) refutes the whole component, where the elimination
of~\cite{m4} refutes one state per round.  The loss: mutual reachability is
quadratic where the elimination is one removing pass, and on the branch where
the SCC decomposition is the whole story it is measurably slower
(Section~\ref{sec:compare}).

\section{The graph layer}\label{sec:graph}

Throughout this section a \emph{graph} is a finite list $V$ of vertices
together with a decidable edge relation $G$; a \emph{fairness family} is a
finite list $f_1,\dots,f_k$ of decidable predicates on vertices.  No formula
occurs anywhere in this section.

For a sublist $W\subseteq V$ write $\outs_W(x):=\{y\in W: G(x,y)\}$.

\begin{definition}
A \emph{trail} in $W$ from $u$ to $w$ is a finite list $z_1,\dots,z_n$ of
vertices of $W$ with $G(u,z_1)$, $G(z_i,z_{i+1})$ and $z_n=w$; the empty trail
runs from $u$ to $u$.  A \emph{circuit at $r$ in $W$ for $f_1,\dots,f_k$} is a
non-empty trail from $r$ to $r$ together with, for each $i$, a vertex on it
satisfying $f_i$.
\end{definition}

Indexing a trail by its endpoint rather than computing the endpoint with a
$\mathrm{last}$ function is what keeps every concatenation below free of
equational bookkeeping; this is the device of~\cite[\S4c]{m4} again.

\begin{definition}[Certificate]\label{def:cert}
The certificates $\mathrm{Cert}(W)$, indexed by the list $W$ of vertices
\emph{still under consideration}, are generated inductively by
\begin{itemize}
\item[(C0)] $W$ is empty;
\item[(C1)] there is $x\in W$ with $\outs_W(x)=\varnothing$, and a certificate
  for $W\setminus\{x\}$;
\item[(C2)] there are $C\subseteq W$ closed under $G$-successors lying in $W$,
  and an $i$ with $f_i(z)=0$ for every $z\in C$, and a certificate for
  $W\setminus C$.
\end{itemize}
\end{definition}

A certificate is a finite object: it is the removal order, with a reason
attached to each step.  It is \emph{not} a decomposition into components ---
no component object, no quotient and no equality of components occurs --- and
(C2) does not require $C$ to be strongly connected, only successor-closed.

\begin{lemma}[Induction principle]\label{lem:certall}
Let $R$ be a property of vertices.  Suppose
\begin{itemize}
\item[(H1)] for every $W$ and every $x\in W$ with $\outs_W(x)=\varnothing$: if
  $R$ holds of every vertex outside $W$, then $R(x)$;
\item[(H2)] for every $W$, every $C\subseteq W$ closed under $G$-successors in
  $W$ and every $i$ with $f_i$ false throughout $C$: if $R$ holds of every
  vertex outside $W$, then $R$ holds of every $z\in C$.
\end{itemize}
Then a certificate for $W$, together with $R$ outside $W$, gives $R$ of every
vertex.
\end{lemma}

\begin{proof}
Induction on the certificate.  In case (C0) nothing is left, so the hypothesis
already gives $R$ everywhere.  In case (C1) apply (H1) and note that a vertex
outside $W\setminus\{x\}$ is outside $W$ or is $x$.  Case (C2) is the same with
(H2) and $C$ in place of $x$.
\end{proof}

Lemma~\ref{lem:certall} is the whole reason the certificate is defined as it
is: (H1) and (H2) are, on the nose, cases (R1) and (R2)
of~\cite[Def.~4.1]{m4}, with the formulas removed.  That is made precise in
Section~\ref{sec:extract}.

\subsection*{Computing one}

\begin{definition}\label{def:reach}
$\Reach_W(v)$ is computed by a \emph{growing} search: start from the empty
accumulator; a round adds every vertex of $W$ that is wanted --- is $v$, or is
a $G$-successor of something already accumulated --- and is not yet there; stop
when a round adds nothing.
\end{definition}

The stopping condition, not a bound, is what gives completeness: when a round
adds nothing the accumulator contains $v$ and is closed under successors in
$W$, and an induction along a trail then shows it contains everything reachable.
There is no fuel and no claim that $|W|$ rounds suffice.  Termination is by the
inductive finiteness of vertices: one cannot go on choosing distinct vertices
from a fixed finite list.  This is the device of~\cite[\S4g]{m4}, applied to the
one computation that note deliberately avoided.

\begin{definition}
$\mathrm{scc}_W(r):=\{w\in\Reach_W(r): r\in\Reach_W(w)\}$, and $r$ is
\emph{terminal} in $W$ when $\Reach_W(r)=\mathrm{scc}_W(r)$, i.e.\ when
everything $r$ reaches reaches $r$ back.
\end{definition}

Both are decidable, and the condensation is never built: a condensation edge is
followed by \emph{moving the representative}, and terminality is one boolean.

\begin{lemma}\label{lem:findterm}
A non-empty $W$ contains a terminal representative.
\end{lemma}

\begin{proof}
While $r$ is not terminal there is $w\in\Reach_W(r)$ with $r\notin\Reach_W(w)$;
move the representative to $w$.  The search terminates because
$\Reach_W(w)$ avoids $r$, so each move gives up one vertex of the ambient list,
and one cannot go on giving up distinct vertices.  The invariant carried is a
single hypothesis --- every vertex of the current reachable set still satisfies
the bar predicate --- and no list of visited representatives is needed; that
acyclicity of the condensation makes the components visited distinct is thereby
not a lemma but the termination of this search.
\end{proof}

\begin{lemma}[One step]\label{lem:block}
Suppose every $x\in W$ has an out-edge in $W$.  Then either $W$ is empty, or
its terminal component is fair and a circuit is built in it, or its terminal
component misses some $f_i$ everywhere.
\end{lemma}

\begin{proof}
Take $r$ terminal by Lemma~\ref{lem:findterm} and $C:=\mathrm{scc}_W(r)$.  If
some $f_i$ is false throughout $C$ we are in the third case.  Otherwise each
$f_i$ has a witness $w_i\in C$.  Since $r$ has an out-edge $x\in W$ and $r$
reaches $x$, terminality sends a trail from $x$ back to $r$, so there is a
non-empty closed trail $c_0$ at $r$.  Now add one \emph{leg} per fairness
predicate: the leg for $f_i$ is $c_0$ followed by a trail $r\to w_i$ inside $C$
followed by a trail $w_i\to r$ inside $C$, both of which exist because $C$ is a
component.  Carrying $c_0$ in front makes the first half non-empty, so $w_i$ is
a vertex the leg \emph{visits} and not merely an endpoint --- which matters when
$w_i=r$.  The concatenation of the legs with $c_0$ is a circuit at $r$.
\end{proof}

\begin{theorem}[Fair-graph certificate]\label{thm:fairgraph}
Given a graph, a fairness family and a vertex $v\in V$, one can compute either
\begin{itemize}
\item[(a)] a lasso from $v$ --- a trail from $v$ to some $r$, and a circuit at
  $r$ --- whose cycle meets every $f_i$, or
\item[(b)] a certificate for $\Reach_V(v)$.
\end{itemize}
\end{theorem}

\begin{proof}
Run the following removing iteration on $W:=\Reach_V(v)$.  Ask whether some
$x\in W$ has $\outs_W(x)=\varnothing$.  If so, recurse on $W\setminus\{x\}$;
a lasso found there is a lasso in $W$, and a certificate found there is
extended by (C1).  If not, apply Lemma~\ref{lem:block}: the empty case gives
(C0), the fair case gives a circuit, and the unfair case gives a $C$ and an $i$,
and we recurse on $W\setminus C$, extending by (C2).  Each step removes at least
one vertex --- for (C2) the representative $r$, which lies in its own component
--- so the iteration terminates.  Finally a circuit at $r$ inside $\Reach_V(v)$
is joined to $v$ by the trail that put $r$ in the reachable set.
\end{proof}

Note that no rank function and no topological sort appear: the iteration
removes, so termination is by shrinking, and the only thing that needed proof
is Lemma~\ref{lem:findterm}.

\subsection*{The two answers are exclusive}

This is the half that replaces the lecture's ``key insight'', and the point is
that it comes out finitary.

\begin{theorem}\label{thm:nocirc}
A certificate for $W$ forbids a circuit inside $W$.
\end{theorem}

\begin{proof}
Induction on the certificate; each step shows the circuit \emph{survives} the
removal, so it is still there when nothing is left, which (C0) forbids.

For (C1): a vertex on a cycle has a successor --- along a trail, every vertex
visited has a successor provided the endpoint has, and for a closed trail the
endpoint is the start, whose successor is the trail's own first step --- so a
sink lies on no cycle, and the cycle avoids $x$.

For (C2): any vertex of a cycle reaches any other --- take the suffix of the
trail to the base point and then the prefix out again --- and a successor-closed
block that a trail starts in contains everything the trail reaches.  So if the
cycle met $C$ it would lie inside $C$, and then its witness for $f_i$ would
satisfy an $f_i$ that is false throughout $C$.  Hence the cycle avoids $C$.
\end{proof}

\begin{corollary}\label{cor:decision}
The negative answer of Theorem~\ref{thm:fairgraph} is a decision: if a fair
lasso from $v$ exists at all, answer (a) is returned.
\end{corollary}

\begin{proof}
A lasso from $v$ has all of its cycle inside $\Reach_V(v)$, by the completeness
of Definition~\ref{def:reach}; so a certificate for $\Reach_V(v)$ would
contradict Theorem~\ref{thm:nocirc}.
\end{proof}

Both branches are inhabited and neither theorem is vacuous: the graph
$p\to q$, $q\to s$, $s\to q$ with the two predicates ``is $q$'', ``is $s$''
exhibits a circuit, and with the single predicate ``is $p$'' exhibits a
certificate, whose removal order is (C2), (C1), (C0) --- the recurrent
component $\{q,s\}$ is refuted first, and only then is $p$ a sink.  That is the
backwards induction in miniature: $p$ is refuted \emph{because} its successors
already were.

\section{Where relevance goes}\label{sec:evpred}

Fix $\varphi$ and let $\Always A_1,\dots,\Always A_m$ enumerate the members of
$\Sub(\varphi)$ of the form $\Always A$, as in~\cite{m4}.  The graph is taken
to be the states with the transition relation of~\cite[Def.~2.3]{m4}.

The informal argument quantifies over ``the fairness sets \emph{relevant} in
$C$'', which is well defined only because a pending eventuality is uniform on a
component: if $\Always A\notin v$ for one $v\in C$ then, by persistence along a
path plus strong connectivity, $\Always A\notin w$ for every $w\in C$.  We do
not need that lemma, because relevance can be put inside the predicate.

\begin{definition}\label{def:evpred}
For each $j$ let
\[
  f_j(z) \;:=\; [\,\Always A_j\in z\,]\ \vee\ [\,A_j\notin z\,],
\]
that is: the $j$-th eventuality is not pending at $z$, or it is fulfilled at
$z$.
\end{definition}

A component in which the eventuality is not pending meets $f_j$ for free, and a
component in which it is pending meets $f_j$ exactly by fulfilling it.  So the
graph layer of Section~\ref{sec:graph} knows only ``$C$ meets every $f$'', and
uniformity on a component is never needed.  The two readings of the same
definition are what Sections~\ref{sec:lasso} and~\ref{sec:extract} use: its
\emph{truth} at a witness is fulfilment, and its \emph{falsity} throughout a
block is exactly the pair of hypotheses of case (R2).

\section{The lasso branch}\label{sec:lasso}

A trail in the state graph \emph{is} a path in the sense of~\cite[\S6]{m4}; the
only work is the cycle-fulfilment property that
Proposition~\ref{prop:lasso5.3} below requires.

\begin{lemma}[Persistence to a visited state]\label{lem:persistto}
Let $\Always A\notin v$ and let $z$ be a state visited by a path
$v\to\cdots\to x$.  Then either the eventuality is fulfilled at some state
visited, or $\Always A\notin z$.
\end{lemma}

\begin{proof}
The induction of~\cite[Lem.~2.8]{m4} propagated along the path, with the target
carried along and the path cut short when it is reached.
\end{proof}

The corresponding statement for the \emph{endpoint} is the one~\cite{m4} needs;
the witness that Theorem~\ref{thm:fairgraph} supplies lies somewhere on the
cycle and not at its base point, which is why the visited-state form is needed
here.

\begin{proposition}\label{prop:lasso5.3}
A fair lasso in the sense of Theorem~\ref{thm:fairgraph}, for the predicates of
Definition~\ref{def:evpred} and inside any set of states contained in
$U_\infty$, is a fulfilling lasso in the sense of~\cite[Prop.~5.3]{m4}.
\end{proposition}

\begin{proof}
The stem and the cycle are paths, and the cycle is non-empty, so only the
cycle-fulfilment property has to be produced: for a state $s$ on the cycle at
which $\Always A_j$ is pending, a state of the cycle at which $A_j$ fails.
Travel along the cycle from $s$ to the base point $r$: by~\cite[Lem.~2.8]{m4}
either $A_j$ fails on the way --- and we are done --- or $\Always A_j\notin r$.
In the latter case go out from $r$ to the witness $w$ that
Theorem~\ref{thm:fairgraph} chose for $f_j$: by Lemma~\ref{lem:persistto},
either $A_j$ fails on the way, or $\Always A_j\notin w$, and then the first
disjunct of $f_j(w)$ is false, so the second holds and $A_j\notin w$.
\end{proof}

So the positive branch feeds the unchanged tail of~\cite{m4}: the lasso is
unwound to a run, Hintikka's lemma evaluates $\neg\varphi$ to $1$ at the start,
and soundness turns any derivation into a contradiction.

\begin{remark}\label{rem:converse}
The converse also holds, and is worth stating because it makes
Proposition~\ref{prop:lasso5.3} unconditional.  A fulfilling lasso is a fair
lasso for the $f_j$: if the eventuality is not pending at the base point, the
base point is its own witness --- and it lies \emph{on} the cycle, the cycle
being non-empty --- and if it is pending, the lasso's own cycle-fulfilment
property supplies a state of the cycle at which $A_j$ fails, which is the other
disjunct.  Hence, by~\cite[Prop.~5.3]{m4} and Corollary~\ref{cor:decision},
Theorem~\ref{thm:fairgraph} run on $U_\infty$ \emph{always} returns answer (a),
and one obtains a second proof of~\cite[Prop.~5.3]{m4} in which no round and no
pigeonhole occurs.
\end{remark}

\section{The refutation branch}\label{sec:extract}

Take $R(v):=\ \proves\neg\chi_v$ in Lemma~\ref{lem:certall}.  Then the
hypothesis carried --- $R$ holds outside $W$ --- is the induction hypothesis of
\cite[Thm.~4.1]{m4} verbatim, and the two hypotheses to be discharged are its
two cases.

\medskip\noindent\textbf{(H1) is case (R1).}
The out-neighbours of $x$ inside $W$ are the $\to$-successors of $x$ inside
$W$, so ``$\outs_W(x)=\varnothing$'' \emph{is} the hypothesis of (R1).  The
proof is~\cite[Thm.~4.1, case (R1)]{m4} unchanged: the step lemma offers all
state-successors, each is refuted by the induction hypothesis, (D1) puts the
refutation under $\Next$ and (Lin) brings the negation out.

\medskip\noindent\textbf{(H2) is case (R2), improved.}
Let $C$ be the block and $f_j$ the predicate it misses.  Then for every $z\in C$
\[
  f_j(z)=0
  \qquad\text{i.e.}\qquad
  \Always A_j\notin z \quad\text{and}\quad A_j\in z ,
\]
which are the two hypotheses (R2) needs, \emph{at every state of $C$} rather
than at one state per round.  Take $B:=\beta_C$; then
\begin{itemize}
\item[(i)] $\proves B\to A_j$, since $A_j\in z$ throughout $C$;
\item[(ii)] $\proves B\to\Next B$, by the step lemma at each $z\in C$, dropping
  the successors outside $W$ by the induction hypothesis and keeping the rest
  because $C$ is closed under successors lying in $W$;
\item[(iii)] $\proves B\to\Always A_j$, by (Ind) through (D5) and (D3).
\end{itemize}
Since every $z\in C$ has $\Always A_j$ pending, $\proves\chi_z\to\neg\Always
A_j$, and $\proves\chi_z\to B$; so $\proves\neg\chi_z$ for \emph{every} $z\in C$
from a single use of (Ind).
\medskip

This is the improvement the SCC form buys.  In~\cite{m4} the greatest fixed
point $W_j$ is recomputed each round and refutes the states pending in it one
round at a time; here the whole component leaves at once, with one invariant.

\begin{remark}
The reader will notice that Lemma~\ref{lem:certall} quantifies (H1) and (H2)
over an arbitrary $W$ and knows nothing about states --- it is the graph
layer's induction principle and has no business knowing.  The formalisation
therefore intersects the block with the ambient state set on the spot; the
intersection is still successor-closed, because the state set is.  It is one
filter and no new theory.
\end{remark}

\section{Two lemmas at the seam}\label{sec:seam}

The two branches want different sets.  The refutation branch is stated against
the set the search explores --- in the formalisation $\Ust(\varphi)$, the
$\to$-closure of the states containing $\varphi$, which replaces $\St(\varphi)$
for the search but not for the derivations --- while the lasso branch delivers
into machinery stated against $U_\infty$.  Two lemmas bridge them, and neither
is in~\cite{m4}.

\begin{lemma}\label{lem:seeded}
Every state of $\Ust(\varphi)$ is reached by a path from a state containing
$\varphi$.
\end{lemma}

\begin{proof}
The closure is computed by the growing search of Definition~\ref{def:reach}
with the seed changed; the same recursion shows that a state a round adds is
either a seed --- and then the empty path does it --- or a successor of
something already accumulated, and then that thing's path grows by one step.
\end{proof}

Lemma~\ref{lem:seeded} says the closure contains \emph{nothing else}, which is
what makes ``a fair component anywhere in $\Ust(\varphi)$'' the right question
to ask: without it the positive branch would exhibit a loop that is a
countermodel to nothing.

\begin{definition}
A set $C$ of states is a \emph{trap} if every $z\in C$ has a
$\to$-successor in $C$, and for every $z\in C$ and every $j$ pending at $z$
there is a path inside $C$ from $z$ to a state fulfilling the $j$-th
eventuality.
\end{definition}

\begin{lemma}\label{lem:trap}
A trap $C\subseteq\St(\varphi)$ contained in $\Ust(\varphi)$ is contained in
$U_\infty$.
\end{lemma}

\begin{proof}
The two clauses of a trap are the negations of (R1) and (R2) read as properties
of a \emph{set}: if $C\subseteq U_k$ then no $z\in C$ is removed by (R1),
having a successor in $C\subseteq U_k$; and none is removed by (R2), since a
path inside $C\subseteq U_k$ to a fulfilling state contradicts the soundness of
the greatest fixed point of~\cite[Rem.~3.2]{m4}.  So $C\subseteq U_{k+1}$, and
the claim follows by induction along the elimination.
\end{proof}

Lemma~\ref{lem:trap} is one direction of the completeness of the elimination:
it says the elimination throws nothing useful away.  The other direction ---
that $U_\infty$ is itself a trap --- is~\cite[Lem.~5.1]{m4}.  Note that the
vertex set of a fair lasso is a trap: every vertex has the next one as a
successor, the base point has the cycle's first step, and a pending eventuality
is fulfilled along the walk to the witness by
Proposition~\ref{prop:lasso5.3}'s argument.

\section{The decision procedure}\label{sec:decide}

\begin{theorem}\label{thm:sccmain}
There is a total computable function taking $\varphi$ to an element of
\[
  \mathrm{Deriv}_{\mathcal{L}}(\varphi)\;+\;\mathrm{LassoCountermodel}(\varphi),
\]
which decides by asking which answer Theorem~\ref{thm:fairgraph} gives on
$\Ust(\neg\varphi)$.
\end{theorem}

\begin{proof}
Run the iteration of Theorem~\ref{thm:fairgraph} on the whole of
$\Ust(\neg\varphi)$ --- not from a single vertex, since the refutation branch
must account for every state at once --- with the predicates of
Definition~\ref{def:evpred}.

If it returns a certificate for $\Ust(\neg\varphi)$, apply
Lemma~\ref{lem:certall} with $R(v):=\ \proves\neg\chi_v$ and
Section~\ref{sec:extract}; the hypothesis outside $\Ust(\neg\varphi)$ is
vacuous.  Every state of $\Ust(\neg\varphi)$ is then refutable, and a state
outside it cannot contain $\neg\varphi$, so $\proves\chi_v\to\varphi$ for all
of them by (L2); conclude by~\cite[Lem.~2.5]{m4} and (L3),(L5).  Observe that
the elimination of~\cite{m4} is not run at all on this branch.

If it returns a fair component at $r$, take by Lemma~\ref{lem:seeded} a state
$s$ containing $\neg\varphi$ and a path $s\to\cdots\to r$; stem and cycle
together are a fair lasso from $s$.  Its vertex set is a trap, so by
Lemma~\ref{lem:trap} it lies in $U_\infty$, and
Proposition~\ref{prop:lasso5.3} turns it into a fulfilling lasso through $s$.
Hintikka's lemma and soundness then give the countermodel exactly as
in~\cite[Thm.~6.1]{m4}.
\end{proof}

The two branches are asymmetric in the same way as in~\cite{m4}: the
certificate branch returns a derivation and never mentions a lasso, and the
fair branch returns a lasso and mentions a derivation only to consume one.
Note the asymmetry in cost as well.  On the certificate branch the elimination
is never computed; on the fair branch it \emph{is} computed, by
Lemma~\ref{lem:trap}, but it is not what decides and it produces no lasso.

\section{Comparison}\label{sec:compare}

The two procedures are in one development, so they can be compared by checked
equation rather than by running them and looking.

\medskip\noindent\textbf{Agreement.}  On the forty-six statements of Figure~11
of~\cite{sgl}, together with
$\Always\Event(p\vee q)\to(\Always\Event p\vee\Always\Event q)$, the two
procedures take the same branch; each of the forty-seven equations
$\mathrm{isProof}(\mathrm{decideSCC}\,X)=\mathrm{isProof}(\mathrm{decide}\,X)$
is checked.  Twenty-two are proofs and twenty-five are countermodels.

\medskip\noindent\textbf{Different lassos.}  For
$\Always(p\vee q)\to(\Always p\vee\Always q)$ the two countermodels differ:
\begin{center}
\begin{tabular}{lccc}
              & start & stem & cycle\\\hline
$\mathrm{decide}$    & $p\wedge q$, neither forever & 3 & 1\\
$\mathrm{decideSCC}$ & $q$, not $p$                 & 2 & 4\\
\end{tabular}
\end{center}
Both are countermodels and each carries its own refutation of derivability.
The cycle of length four is not an accident: the construction of
Lemma~\ref{lem:block} prepends one leg per fairness predicate to the bare loop,
and $\Sub$ of the negation has three eventualities.

\medskip\noindent\textbf{Cost.}  Measured cold, with the import baseline
subtracted:
\begin{center}
\begin{tabular}{lcc}
   & $\mathrm{decide}$ & $\mathrm{decideSCC}$\\\hline
$\Always(p\vee q)\to(\Always p\vee\Always q)$ & $0.19$ s & $0.18$ s\\
$\Always\Event(p\vee q)\to(\Always\Event p\vee\Always\Event q)$
                                              & $0.78$ s & $1.24$ s\\
\end{tabular}
\end{center}
The first is decided on the countermodel branch and the second on the proof
branch.  The prediction that the SCC form is slower --- mutual reachability
being quadratic in the component where the greatest fixed point is one removing
pass --- is confirmed on the proof branch, where it is $1.6$ times slower, and
refuted on the countermodel branch, where the two are level: there the SCC
procedure pays for both the graph search and Lemma~\ref{lem:trap}, but saves
the round construction and the boundary search of~\cite[Prop.~5.3]{m4}, and on
that formula the two cancel.

\section{Remarks}

\begin{remark}[What the phase buys, and what it does not]
It does not buy constructivity: \cite{m4} is already constructive, and its
greatest fixed point is arguably the SCC decomposition performed without
decomposing --- the union of the bad recurrent regions, computed by a removing
iteration in one pass.  It does not buy speed.  What it buys is the factoring
of Section~\ref{sec:graph}, an explanation of \emph{why} (Ind) appears at all,
and the refutation of a whole component by one invariant.
\end{remark}

\begin{remark}[Absences]
No infinite path is constructed; nothing is coinductive; there is no pigeonhole
lemma; there is no rank function and no topological sort; there is no component
object, no quotient and no equality of components; and the condensation is
never built.  Definition~\ref{def:reach} is the one computation
\cite[\S4c]{m4} declined to do, and it is done without cardinality theory.
\end{remark}

\begin{remark}[A worked example, and a correction]\label{rem:lecture}
The appendix of the lecture notes~\cite{lec} analyses
$\Event a\wedge\Event b\to\Event(a\wedge b)$ by hand: it lists the locally
consistent states with $\Event(a\wedge b)=0$, their out-edges, and a
countermodel path.  The graph agrees with the one computed here exactly ---
eight states with the same valuations, twenty-four edges, the same out-degrees
--- although the two are obtained differently: the hand analysis \emph{assumes}
$\Event(a\wedge b)=0$ to cut its table down, whereas here the eight states are
the closure of the states containing the negation and $\Event(a\wedge b)=0$
falls out.

The countermodel path given there, $s_3\to s_5\to s_7\to s_7\to\cdots$ looping
at $s_7$, is however \emph{not fair}: at $s_7$ one has $\Event b=1$ and $b=0$,
so $s_7$ lies outside the very fairness set $F_{\Event b}$ that the notes
define, and the loop promises $b$ forever without delivering it.  In the
notation of Definition~\ref{def:evpred} this is the checked equation
$f_{\Event b}(s_7)=0$.  The path presumably intended is
$s_3\to s_5\to s_7\to s_5\to\cdots$, alternating so that $b$ recurs at $s_5$ and
$a$ at $s_7$; both edges exist.  The procedure of
Theorem~\ref{thm:sccmain} returns instead $s_7\to s_4\to s_0\to s_0\to\cdots$:
$a$ at stage $0$, $b$ at stage $1$, and everything false from stage $2$ on.  It
ends in a state where nothing is pending at all, which the hand analysis never
considers because it is looking for a loop that \emph{keeps} the eventualities
alive.
\end{remark}

\begin{remark}[Formalisation]
Everything above has been checked in Agda --- \texttt{--safe},
\texttt{--without-K}, \texttt{--exact-split}, no library, no postulates, no
unfilled goals, no termination pragmas, and no \texttt{with} --- beside the
development of~\cite{m4} and sharing its step lemma, its case (R1), its notion
of run and its soundness.  The new modules are
\begin{center}
\begin{tabular}{llrr}
\texttt{LTL.Graph}      & Section~\ref{sec:graph}     & 1225 & 0.60 s\\
\texttt{LTL.SccLasso}   & Sections~\ref{sec:evpred}, \ref{sec:lasso}
                                                      &  372 & 0.34 s\\
\texttt{LTL.SccExtract} & Section~\ref{sec:extract}   &  380 & 0.32 s\\
\texttt{LTL.SccTrap}    & Section~\ref{sec:seam}      &  289 & 0.21 s\\
\texttt{LTL.DecideSCC}  & Section~\ref{sec:decide}    &  256 & 0.30 s\\
\texttt{LTL.Compare}    & Section~\ref{sec:compare}   &  408 & 20.7 s\\
\texttt{LTL.Lecture}    & Remark~\ref{rem:lecture}    &  184 & 10.2 s\\
\end{tabular}
\end{center}
lines and cold check time; the last two are benchmarks, being forty-seven
decision problems and one worked example respectively.  \texttt{LTL.Graph}
contains no occurrence of the type of formulas.

Three places where formalising changed the argument rather than transcribing
it.  Lemma~\ref{lem:persistto}, which the informal argument does not
distinguish from persistence to the endpoint, and which is needed because the
witness sits on the cycle rather than at its base point.
Lemma~\ref{lem:seeded}, which the closure construction of~\cite{m4} did not
have --- it proved that the closure contains the seeds and is closed, but not
that it contains nothing else --- and without which the fair branch exhibits a
loop that is a countermodel to nothing.  And Lemma~\ref{lem:trap}, which the
informal accounts do not mention at all because they do not have two state
sets to bridge; it is the one place where the SCC procedure must still compute
the elimination of~\cite{m4}.
\end{remark}

\begin{remark}
The construction is stated for the $\Always$-only syntax of~\cite{m4}, where
persistence is (T3) and is already proved.  Extending it to \emph{until} faces
exactly the obligation recorded at the end of~\cite{m4}, and one more: the
fairness predicate of Definition~\ref{def:evpred} would become
$[\,A\,\mathcal{U}\,B\notin z\,]\vee[\,B\in z\,]$, and the argument of
Proposition~\ref{prop:lasso5.3} would have to carry the side condition that $A$
holds along the way.
\end{remark}

\begin{thebibliography}{9}
\bibitem{m4} \emph{Proof extraction from a failed LTL tableau: a uniform
loop-invariant construction}, companion note.
\bibitem{lec} \emph{A Complete Decision Algorithm for Linear Temporal Logic},
lecture notes.
\bibitem{sgl} P.~H.~Schmitt and J.~Goubault-Larrecq, \emph{Linear Temporal
Logic}, Figure~11.
\end{thebibliography}

\end{document}
